\documentclass{article}
\usepackage{amsmath, amssymb, amsthm}
\usepackage{hyperref}
\usepackage{geometry}
\geometry{margin=1in}

\title{Erd\H{o}s Problem \#887}
\date{Last updated: 2025-08-31}
\author{Source: erdosproblems.com}

\begin{document}
\maketitle

\noindent\textbf{Status:} OPEN \quad \textbf{No prize}

\noindent\textbf{Tags:} number theory, divisors

\medskip
\noindent\textbf{Problem Statement:}

Is there an absolute constant $K$ such that, for every $C>0$, if $n$ is sufficiently large then $n$ has at most $K$ divisors in $(n^{1/2},n^{1/2}+C n^{1/4})$.

\bigskip
\noindent\textbf{Remarks:}

A question of Erd\H{o}s and Rosenfeld \cite{ErRo97}, who proved that there are infinitely many $n$ with $4$ divisors in $(n^{1/2},n^{1/2}+n^{1/4})$, and ask whether $4$ is best possible here. They also proved that, if $n$ is large enough depending on $C$, then $n$ has at most $1+C^2$ divisors in $(n^{1/2},n^{1/2}+C n^{1/4})$. 

Chan \cite{Ch14} has resolved this when $n$ is a square, proving that if $n$ is a square then $n$ has at most $5$ divisors in\[[ n^{1/2}-n^{1/4}(\log n)^{1/7}, n^{1/2}+n^{1/4}(\log n)^{1/7}].\]In \cite{Ch15} Chan further proves that if $n=(N-a)(N-b)$ for some $0\leq a\leq b\leq \exp((\log n)^{2/7})$ then $n$ has at most $18$ divisors in\[[ n^{1/2}-n^{1/4}(\log n)^{1/14}, n^{1/2}+n^{1/4}(\log n)^{1/14}].\]See also [886].

\bigskip
\noindent\textbf{References:}

[Ch14] Chan, Tsz Ho, Factors of a perfect square. Acta Arith. (2014), 141--143.
 
 [Ch15] Chan, Tsz Ho, Factors of almost squares and lattice points on circles. Int. J. Number Theory (2015), 1701--1708.
 
 [ErRo97] Erd\H{o}s, Paul and Rosenfeld, Moshe, The factor-difference set of integers. Acta Arith. (1997), 353--359.

\bigskip
\noindent\small{Source: \url{https://www.erdosproblems.com/887}}
\end{document}
